\begin{thebibliography}{99}
\addcontentsline{toc}{chapter}{Bibliography}
\bibitem{adleman} Leonard M. Adleman. Molecular computation of solutions to combinatorial problems. Science 266 (1994), 1021--1024. \href{https://computingbiology.github.io/docs/adleman1994.pdf}{Source link}.

\bibitem{dna} Bruce Alberts et al. The Structure and Function of DNA. Molecular Biology of the Cell, 4th edition (2002). \href{https://www.ncbi.nlm.nih.gov/books/NBK26821/}{Source link}.

\bibitem{ligase} New England Biolabs. T4 DNA Ligase, product description. \href{https://www.neb.com/en-us/products/m0202-t4-dna-ligase}{Source link}.

\bibitem{pcr} National Human Genome Research Institute. Polymerase Chain Reaction Fact Sheet. \href{https://www.genome.gov/about-genomics/fact-sheets/Polymerase-Chain-Reaction-Fact-Sheet}{Source link}.

\bibitem{gel} National Human Genome Research Institute. Electrophoresis. \href{https://www.genome.gov/genetics-glossary/Electrophoresis}{Source link}.

\bibitem{hamilton} Frank Pfenning. Lecture 36: The Hamiltonian Path Problem. Carnegie Mellon University, 15-453. \href{https://www.cs.cmu.edu/~fp/courses/flac/lectures/lecture36.html}{Source link}.

\bibitem{si} Bureau International des Poids et Mesures. SI base unit: mole. \href{https://www.bipm.org/en/si-base-units/mole}{Source link}.

\bibitem{benenson} Yaakov Benenson et al. Programmable and autonomous computing machine made of biomolecules. Nature 414 (2001), 430--434. \href{https://www.nature.com/articles/35106533}{Source link}.

\bibitem{crn} David Soloveichik, Georg Seelig, and Erik Winfree. DNA as a universal substrate for chemical kinetics. PNAS 107 (2010), 5393--5398. \href{https://doi.org/10.1073/pnas.0909380107}{Source link}.

\bibitem{circuits} Lulu Qian and Erik Winfree. Scaling up digital circuit computation with DNA strand displacement cascades. Science 332 (2011), 1196--1201. \href{https://www.qianlab.caltech.edu/seesaw_digital_circuits2011.pdf}{Source link}.

\bibitem{assembly} Erik Winfree, Furong Liu, Lisa A. Wenzler, and Nadrian C. Seeman. Design and self-assembly of two-dimensional DNA crystals. Nature 394 (1998), 539--544. \href{https://www.nature.com/articles/28998}{Source link}.
\end{thebibliography}
